\scanpage{69}
\mspage{65}
% source lines 1--3
\[
\begin{aligned}
\alpha_*(\operatorname{supp}^{\unsure{\circ|\bullet}}_E(A))&\overset{?}{\subset}\operatorname{supp}^{\circ}_{E'}(\bar\alpha(A))\\
\Vert\qquad&\\
\operatorname{supp}^{\circ}_{E'}(\bar\alpha(\operatorname{supp}^{\circ}_E(A))&
\end{aligned}
\]
ou encore de ceci:
\[
\left\{\begin{aligned}
&\text{Si }B\subset\operatorname{supp}^{\circ}_E(A)\quad\text{i.e.}\quad\forall x\in E,\ x\mid\circ\mid A\Rightarrow x\mid\circ\mid B\\
&\text{peut-on en conclure que}\\
&?\left\{\begin{aligned}
\bar\alpha(B)&\subset\deleted{\unsure{\bar\alpha(\operatorname{supp}}}\operatorname{supp}^{\circ}_{E'}(\bar\alpha(A))\quad\text{i.e.}\quad\forall x'\in E'\\
&x'\mid\circ\mid\bar\alpha(A)\Rightarrow x'\mid\circ\mid\bar\alpha(B
\end{aligned}\right.\\
&\hfill ?
\end{aligned}\right.
\]
% source lines 8--12
Pour l'affirmer, il faudrait une corr. en sens inverse
\diagram{\begin{tikzcd}[ampersand replacement=\&]
\beta:E' \arrow[r,no head] \& \unsure{\mathcal{P}|\mathscr{P}}(E)
\end{tikzcd}}
telle que l'on ait, $\forall x\in E,\ x'\in E'$
\[
x\mid\circ\mid\beta(x')\Longleftrightarrow\alpha(x)\mid\circ\mid x',\tag*{$(*)$}
\]
\inserted{d'où, pour $A\subset E,\ B'\subset E'$
\[
A\mid\circ\mid\bar\beta(\deleted{A'}\unsure{B'|A'})\Longleftrightarrow\bar\alpha(A)\mid\circ\mid\deleted{A'}B'\tag*{$(**)$}
\]}
% source lines 13--17
Dans ce cas, l'implication ? ci-dessus se \deleted{\unsure{lit}} reformule comme
\[
\deleted{\unsure{x'|x}\mid\circ\mid\beta(\unsure{x'|x})\overset{?}{\Longrightarrow}}
\]
\[
\beta(x')\mid\circ\mid A\overset{?}{\Longrightarrow}\beta(x')\mid\circ\mid B
\]
% margin note written slanted
\margin{On n'a \unsure{pas} utilisé ici $x\mid\circ\mid y\Rightarrow\alpha(x)\unsure{\mid\circ\mid}\alpha(y)$!}
qui résulte bien de $B\subset\operatorname{supp}^{\circ}_E(A)$

% source lines 18--23
Ici encore, si une telle corr. $\beta$ existe, il doit y en avoir une optimale: soit en effet
\[
\gamma(x')\unsure{\equiv|=}\{x\in E\mid\alpha(x)\mid\circ\mid x'\},
\]
\margin{cette fois \unsure{définie|défini} en termes de $\alpha$ seul}
alors la propriété $(*)$ sur $\beta$ se lit
\[
x\mid\circ\mid\beta(x')\Longleftrightarrow x\in\gamma(x')\qquad\text{i.e.}
\]
\[
\gamma(x')=\operatorname{cosupp}^{\circ}(\beta(x'))\qquad\deleted{\text{\unsure{on aura}}}
\]
% archive sheet number: 68
